\subsection*{Shared external-input catalogue}
The five central interfaces are implemented and uninstantiated.\par
Their source files are \texttt{External\allowbreak{}Inputs.lean} and \texttt{Orbit\allowbreak{}Inputs.lean}. KM-U, KM-R and KM-M remain planned interfaces. A catalogue entry is not a Lean axiom; the exact actual types and dependency information are in \texttt{docs/ASSUMPTIONS.md}.

\ifDetailed
\textbf{SS: quoted polynomial-phase bound.} Shao--Sun, \emph{Linear equations in Piatetski--Shapiro primes}, arXiv:2605.18676v1, Theorem 2.1 and Definition 4.5. The source states the coefficient-uniform polynomial-phase consequence directly. \texttt{ShaoSunPhaseInput} quotes that consequence with a saving uniform throughout the degree-dependent interval of $\gamma$. Restricting real cutoffs to integers and extending $X\geq3$ to the two smaller positive cutoffs by enlarging a $\gamma$-dependent constant are included in the external translation. The general nilsequence/circle construction is optional unformalized exposition, not an additional required bridge for this quoted consequence. Both main branches depend on SS.

\textbf{TV: finite polynomial estimate.} Ter\"av\"ainen, \emph{Pointwise convergence of ergodic averages with M\"obius weight}, arXiv:2401.03174v2, Theorem 4.1. Its normalization is $M^{-d-1}$, its support scale is $M^d$, and its control is $u^{d+1}$. \texttt{FiniteVNInput} uses integer cutoffs, positive degrees and a universal bound on every phase correlation. The supremum/sign/monotonicity equivalence with the published norm is the unformalized formulation bridge TV-B. The constant depends on the fixed polynomial family. Proposition 2.3 depends on TV; the pointwise theorem instead quotes OR independently and has no TV premise.

\textbf{KM-U, KM-R and KM-M: planned KMTT inputs.} Krause--Mousavi--Tao--Ter\"av\"ainen, arXiv:2608.19525v1, Proposition 5.2, (5.9)--(5.12), gives weighted-to-unweighted discrepancy. Lemma 5.3, (5.25)--(5.28), gives weak regularity with a retained small-scale alternative. Its coordinate resolutions are $Q^{-O(1)}\varepsilon^{O(1)}M_0^{\nu_i}$, where $\nu_0=\deg P_1$ and $\nu_i=\deg P_i$ for $i\geq1$; the common modulus is $Q^b$. Compatibility requires the enlargements (5.35)--(5.36). The separate mass comparisons are (5.2) and (5.15), retaining any exceptional-zero bias. The local-factor density deduction and quantitative iteration remain original obligations. Only recurrence uses these planned inputs. Precise provisional transcriptions are in \texttt{docs/KMTT\_INPUTS\_WORKING.md}.

\textbf{OR: orbit estimate.} Mousavi, arXiv:2609.01847v1, Proposition 1.4. \texttt{OrbitVNInput} bounds the $L^1$ norm by the scaled finite expression for weights with size and universal phase-correlation bounds. Integer shifts, essential bounds, integrability and the every-phase formulation are included in the unformalized source translation OR-B. The source proof uses duality and orbit averaging; that earlier proof is quoted, not reproved from TV. Only the pointwise branch depends on OR.

\textbf{MW: ordinary weighted pointwise convergence.} Mirek--Wan--Wright, arXiv:2505.15549v2, Theorem 1.1(ii), proves a.e. convergence for finite $1<q_i<\infty$ with $\sum_iq_i^{-1}\leq1$. \texttt{OrdinaryPrimePointwiseInput} uses bounded probability-space inputs and the shifted family $P_i(n-1)$. The finite-$L^q$ embedding and shift specialization are the unformalized translation MW-B. Only Theorem 1.2 depends on MW.

\textbf{PNT: ordinary weighted mass.} \texttt{PrimeNumberTheoremInput} states $N^{-1}\sum_{n\leq N}\Lambda(n)\to1$. Tao's \emph{246B Notes 4: The Riemann zeta function and the prime number theorem}, 12 February 2021, Theorem 6, states this von Mangoldt form. Integer-scale specialization is included in the quoted input. This is not silently inferred from an unrelated prime-counting normalization. Higher-prime-power removal is proved internally. The optional no-Siegel footnote is not used; Green's arXiv:2610.00311v2 Theorem 1.1 is not an additional project premise.

\else
\textbf{SS} is the directly published polynomial-phase consequence of Shao--Sun, arXiv:2605.18676v1, Theorem 2.1 and Definition 4.5. \textbf{TV} is Ter\"av\"ainen, arXiv:2401.03174v2, Theorem 4.1. The universal-phase/supremum bridge TV-B remains external. SS and TV give Proposition 2.3; the pointwise theorem does not require TV.

\textbf{KM-U} is KMTT, arXiv:2608.19525v1, Proposition 5.2; \textbf{KM-R} is Lemma 5.3 with its small-scale alternative and coordinate-dependent resolutions; compatibility uses (5.35)--(5.36). \textbf{KM-M} records (5.2) and (5.15). These interfaces and the original local-factor deduction and quantitative iteration remain open. Only recurrence uses KM.

\textbf{OR} is Mousavi, arXiv:2609.01847v1, Proposition 1.4, via \texttt{OrbitVNInput}. \textbf{MW} is Mirek--Wan--Wright, arXiv:2505.15549v2, Theorem 1.1(ii), via \texttt{OrdinaryPrimePointwiseInput}. Their source translations OR-B and MW-B remain external. \textbf{PNT} is the ordinary von Mangoldt mass limit, directly stated in Tao's 246B Notes 4, Theorem 6. Theorem 1.2 has exactly SS, OR, PNT and MW as its substantial external premises. The optional no-Siegel footnote is not used.
\fi

Primary sources were read in the cited versions. Provenance verification does not formalize an external theorem or discharge its source-formulation bridges. See the repository assumption catalogue and reference review for the actual interfaces.
